\section{Case III: a condensed-phase WCA dimer}
\label{sec:case_wca}

The third system is many-body and is where ABF is most starved: a dimer that must
push a dense solvent aside to change its bond length. It tests whether the
two-dimensional findings survive in a setting with thousands of degrees of freedom
and a nonlinear collective variable.

\subsection{Model and reference}
A dimer of two tagged particles sits in a periodic two-dimensional bath of
$n_{\mathrm{dim}}^2-2$ Weeks--Chandler--Andersen (WCA) solvent particles
($n_{\mathrm{dim}}=10$, lattice spacing $a$, box $L=n_{\mathrm{dim}}\,a$). All
particles interact through the purely repulsive WCA
potential~\citep{WeeksChandlerAndersen1971} (Lennard--Jones cut and shifted at its
minimum $r_0=2^{1/6}\sigma$). The two dimer particles additionally feel a symmetric
double-well bond $V_{\mathrm{dim}}(r)=h\,(1-((r-r_0-w)/w)^2)^2$ with barrier
height $h=2$ and width $w=2$. The reaction coordinate is the scaled bond length
\begin{equation}
  z(q) = \frac{|q_1-q_2|-r_0}{2w},
  \label{eq:wca_cv}
\end{equation}
a nonlinear collective variable of the full configuration $q\in\R^{2n}$
(\cref{rem:manybody}): $z\approx0$ is the \emph{compact} state, $z\approx1$ the
\emph{stretched} state, and $z\in[0.25,0.75]$ the \emph{transition} region. The
free-energy barrier near $z=0.5$ is partly entropic / crowding-driven, since the
dimer must displace solvent to cross it. Dynamics are overdamped Langevin at
$\beta=1$, $\mathrm{d}t=0.002$, with the minimum-image convention; the default
system is $a=1.5$, $\sigma=\epsilon=1$. The local mean force along $z$ is the
constrained projection of the physical forces, including the geometric
contribution with the sign convention of the estimator (\cref{rem:manybody}). The
reference $F(z)$ is computed once per system $a$ by constrained-MD thermodynamic
integration~\citep{Carter1989,DenOtterBriels1998}; it is a high-accuracy numerical
reference used for evaluation only, not an analytic exact solution. Errors are
taken on the interior window
$z\in[-0.1,1.1]$ with the additive constant aligned.

\subsection{Methods and design}
As before the practical comparison is \textbf{ABF only} versus \textbf{ABF+FR
estimated}~\eqref{eq:q_est}, with \textbf{FR uniform} and \textbf{FR oracle}
diagnostics; here the oracle target is built from the TI reference and is
explicitly \emph{not} deployable (\cref{rem:oracle}). A no-leakage assertion
ensures only the oracle method ever receives the reference, and a clone copies the
\emph{entire} WCA configuration --- not just $z$ --- so cloned solvent
arrangements travel with the bond length. Replica lineage is tracked to report the
ancestor effective sample size $\mathrm{ESS}=1/\sum_a w_a^2$. The tuned FR
configuration is given in \cref{tab:wca_design}; the rate is named
\texttt{fr\_rate} and the onset \texttt{fr\_start\_steps} (\cref{sec:fr_notation}).

\input{tables/wca_design.tex}

\subsection{Results}
The main comparison (\cref{tab:wca_main}; 10 seeds, $N=1024$, $250\,000$ steps)
shows the practical method \textbf{reliably beating ABF}: the tuned FR cuts the
median final $\Ltwo(F)$ from $\WCAabfLtwoF$ to $\WCAtunedLtwoF$
(\,$+\WCAtunedGainPct\%$\,), winning all $\WCAtunedWin/10$ matched seeds. As in the
entropic-bottleneck case, the uniform target ($+\WCAuniGainPct\%$) and the oracle
target ($+\WCAoracleGainPct\%$) are \emph{essentially as good as} the estimated
target --- the online estimated target already extracts nearly all the available
benefit, and there is no headroom to the oracle. The stronger estimated-target row
in the same table attains the lowest final free-energy error, but with a much
smaller ancestor ESS; we therefore use the tuned row as the practical headline and
the strong row as the accuracy--diversity contrast. We are explicit that the
headline gain is in the \emph{free energy}: the mean-force error $\Ltwo(F')$
improves only modestly ($\approx22\%$), so we do not claim a large mean-force
improvement.

\input{tables/wca_main.tex}

\begin{figure}[t]
  \centering
  \includegraphics[width=0.92\textwidth]{fig_wca_01_convergence.png}
  \caption{WCA dimer, convergence (median and inter-quartile range over seeds).
    ABF and FR are identical until FR activates; from then on FR is consistently
    below ABF on $\Ltwo(F)$ for the rest of the run. The cumulative replacement
    count (right) confirms the birth--death is active but bounded.}
  \label{fig:wca_convergence}
\end{figure}

\begin{figure}[t]
  \centering
  \includegraphics[width=0.92\textwidth]{fig_wca_02_profiles.png}
  \caption{Final profiles (median over seeds) versus the TI reference (dashed):
    (a) free energy $\Fhat(z)$, (b) mean force $\Fphat(z)$, (c) reaction-coordinate
    marginal. FR recovers the reference free energy more accurately than ABF,
    chiefly in and around the transition region.}
  \label{fig:wca_profiles}
\end{figure}

\subsection{Mechanism: support reallocation with a diversity cost}
The diagnostics support the same broad mechanism as the entropic bottleneck rather
than a stationary-coverage reshaping one. Final reaction-coordinate region
fractions and the kernel-weighted ABF denominator in the transition window are
nearly identical across ABF, FR estimated, FR uniform and FR oracle: FR does
\emph{not} converge to a different stationary $z$-distribution. Instead the gain
appears immediately after FR activates and is sustained for the rest of the run.
\Cref{fig:wca_mechanism} shows the interpretation: birth--death acts as a
\emph{support-reallocation resampler for the ABF estimator}. Replicas stuck in
over-represented configurations are replaced by whole-configuration clones of
under-represented ones, improving weighted force-sample support in important
$z$-bins even when the marginal looks unchanged. A clone is not an independent
solvent arrangement at birth; the point is that useful whole configurations are
moved into starved reaction-coordinate regions and then decorrelated by subsequent
Langevin dynamics. This explains why target shape matters little here, consistent
with $\text{uniform}\approx\text{estimated}\approx\text{oracle}$.

\begin{figure}[t]
  \centering
  \includegraphics[width=0.92\textwidth]{fig_wca_04_mechanism.png}
  \caption{Mechanism. (a) The kernel-weighted ABF denominator
    $N_{\mathrm{eff}}(z)$ is raised by FR in the transition window; (b) per-bin
    mean-force error is reduced there; (c) birth (clone) and death locations
    concentrate around the over/under-sampled regions; FR reshapes the
    \emph{sampling support}, not the stationary marginal.}
  \label{fig:wca_mechanism}
\end{figure}

\subsection{Failure boundary}
Unlike the entropic bottleneck, the WCA dimer exhibits a clear failure boundary in
the FR rate (\cref{fig:wca_failure}). The rate ladder (a shorter $150$k-step,
$4$-seed stage) maps a clean safe regime: too gentle
(\texttt{fr\_rate}$\le0.05$) gives a small gain ($+19$--$29\%$); the sweet spot
(\texttt{fr\_rate}$\approx0.1$--$0.2$) gives the maximal $+58$--$64\%$ with healthy
ancestor ESS; the gain then \emph{erodes} sharply
(\texttt{fr\_rate}$=0.5$ falls to $+12\%$ with the error already climbing back
toward the ABF baseline) and finally \emph{reverses} for
\texttt{fr\_rate}$\ge1$, where the median final $\Ltwo(F)$ is worse than ABF
($-38\%$ at rate $1$, $-68\%$ at rate $2$) and no seed wins, as the ancestor ESS
collapses from $\sim\!337$ to $\sim\!20$. The reversal arrives even earlier under a
longer budget: the production aggressive configuration (\texttt{fr\_rate}$=0.5$ at
$250$k steps, $10$ seeds) is already net-harmful ($\WCAaggrGainPct\%$, winning only
$1/10$ seeds, ancestor ESS $\sim\!22$), because the extra steps give diversity loss
more time to dominate. Over-aggressive birth--death destroys the replica diversity
the mean-force estimator depends on, so the estimator degrades faster than coverage
improves. This is the predicted failure mode; the analytic $y$-channel of the
entropic bottleneck masks it (\cref{sec:case_eb}) because cloned walkers
re-equilibrate instantly, whereas a dense solvent does not.

\begin{figure}[t]
  \centering
  \includegraphics[width=0.92\textwidth]{fig_wca_05_failure.png}
  \caption{FR-rate ladder ($150$k steps, $4$ seeds). Final $\Ltwo(F)$ (left) is
    minimised in the $0.1$--$0.2$ sweet spot, climbs back toward the ABF baseline
    (dashed) by \texttt{fr\_rate}$=0.5$, and rises above it for
    \texttt{fr\_rate}$\ge1$; ancestor ESS (right) collapses as the rate grows,
    pinpointing diversity loss as the failure mechanism.}
  \label{fig:wca_failure}
\end{figure}

\subsection{Difficulty dependence}
Varying one difficulty axis at a time (\cref{fig:wca_difficulty}) shows the gain
is robust and, importantly, \emph{grows with difficulty}. With \emph{budget}, FR
wins at every $n_{\mathrm{steps}}$ and the relative advantage grows rather than
washes out ($+\WCAbudgetLoGain\%$ at $50$k to $+\WCAbudgetHiGain\%$ at $400$k):
ABF's error plateaus around $0.08$ after $\sim150$k while FR keeps descending, so
FR is not merely an early-time accelerator that ABF later catches. With
\emph{replica count}, the gain is roughly flat ($+51$--$59\%$ from $N=256$ to
$2048$) --- FR helps even at $N=256$. With \emph{solvent crowding}, the densest
bath ($a=1.35$) makes the entropic barrier severe and ABF $\sim4\times$ worse in
absolute error; FR still wins $5/5$ and delivers by far the largest absolute error
reduction ($\WCAdenseAbfErr\to\WCAdenseFrErr$). In the units a practitioner cares
about --- absolute free-energy error removed --- the crowded regime is where FR
helps most.

\begin{figure}[t]
  \centering
  \includegraphics[width=0.95\textwidth]{fig_wca_06_difficulty.png}
  \caption{Difficulty dependence (matched-seed gain over ABF). (a) Budget: the
    relative gain grows with $n_{\mathrm{steps}}$ as ABF plateaus and FR keeps
    improving. (b) Replica count: gain is roughly flat. (c) Solvent crowding:
    smaller $a$ (denser) yields the largest absolute error reduction. FR wins
    $5/5$ seeds in every cell.}
  \label{fig:wca_difficulty}
\end{figure}


\subsection{WCA phase diagram: when does marginal Fisher--Rao help?}
\label{sec:wca_phase_diagram}
% Auto-generated by scripts/make_phase_report_assets.py. Do not edit by hand.
% Macros for the WCA phase-diagram subsection (prefix WCAPD).
\newcommand{\WCAPDNCells}{14}
\newcommand{\WCAPDNPlaneCells}{12}
\newcommand{\WCAPDModalM}{100}
\newcommand{\WCAPDNSeeds}{4}
\newcommand{\WCAPDNSteps}{120000}
\newcommand{\WCAPDMeanGain}{-68.70}
\newcommand{\WCAPDMedianGain}{-17.78}
\newcommand{\WCAPDNWinCells}{5}
\newcommand{\WCAPDNLoseCells}{7}
\newcommand{\WCAPDMaxGain}{51.27}
\newcommand{\WCAPDMaxGainCell}{$\beta{=}1,\,h{=}4$}
\newcommand{\WCAPDMaxGainR}{2.05}
\newcommand{\WCAPDMinGain}{-489}
\newcommand{\WCAPDMinGainCell}{$\beta{=}4,\,h{=}1$}
\newcommand{\WCAPDAnchorGain}{47.96}
\newcommand{\WCAPDAnchorR}{1.94}
\newcommand{\WCAPDAnchorAbfErr}{0.0918}
\newcommand{\WCAPDAnchorFrErr}{0.0474}
\newcommand{\WCAPDEasiestCell}{$\beta{=}1,\,h{=}1$}
\newcommand{\WCAPDEasiestBh}{1.00}
\newcommand{\WCAPDEasiestAbfErr}{0.101}
\newcommand{\WCAPDHardestCell}{$\beta{=}4,\,h{=}6$}
\newcommand{\WCAPDHardestBh}{24.00}
\newcommand{\WCAPDHardestAbfErr}{0.0249}
\newcommand{\WCAPDEasiestGain}{49.87}
\newcommand{\WCAPDHardestFrErr}{0.0414}
\newcommand{\WCAPDHardestGain}{-69.10}
\newcommand{\WCAPDCorrGainBh}{-0.121}
\newcommand{\WCAPDMeanAbsUniMinusEst}{10.77}
\newcommand{\WCAPDMeanAbsOraMinusEst}{83.81}
\newcommand{\WCAPDMinEss}{149}
\newcommand{\WCAPDMinFinalEss}{149}
\newcommand{\WCAPDMaxAncFrac}{0.0273}
\newcommand{\WCAPDMeanFinalEss}{213}
\newcommand{\WCAPDHasMassAxis}{1}
\newcommand{\WCAPDMassList}{49,100,196}


A single WCA setting is not enough to tell whether the gain above is a generic
property of the dimer or a response to sample starvation. We therefore added an
additive, command-line phase-diagram study that sweeps physical parameters while
keeping the existing production result intact. In the production phase diagram reported here we vary
the inverse temperature $\beta\in\{1,2,4\}$ and dimer barrier height
$h\in\{1,2,4,6\}$ at fixed width $w=2$, lattice spacing $a=1.5$, and
$M=n_{\mathrm{dim}}^2=100$ physical particles; we also add a targeted
physical-size slice $M\in\{49,100,196\}$ at $(\beta,h)=(2,4)$. This $M$ is the number of physical particles
in one WCA system; it is distinct from the $N=1024$ ABF/FR
replicas used to estimate the mean force. Each cell uses $\WCAPDNSeeds$ matched
seeds and $\WCAPDNSteps$ steps. The deployable comparison is ABF versus
mFR--ABF with the online estimated target; the uniform target is an ablation, and
the oracle target is included as a diagnostic positive control in the table but is
not a usable method because it reads the TI reference.

\begin{table}[t]
\centering
\small
\caption{WCA phase diagram over $(\beta,h)$ at $M=100$ physical particles
($N=1024$ replicas, $120{,}000$ steps).
$R=L^2(F)_{\rm ABF}/L^2(F)_{\rm mFR\,est}$; gain is the matched-seed median
percentage reduction in final $L^2(F)$. The oracle column is a diagnostic that
uses the TI reference and is \emph{not} deployable.}
\label{tab:wca_phase_main}
\begin{tabular}{rrr rrrr rr rr}
\toprule
$\beta$ & $h$ & $\beta h$ & ABF & mFR est & mFR uni & mFR ora$^{*}$ & $R$ & gain (\%) & wins & ESS \\
\midrule
1 & 1 & 1 & 0.1006 & 0.0503 & 0.0518 & 0.0521 & 2.00 & $+49.9$ & 4/4 & 171 \\
1 & 2 & 2 & 0.0918 & 0.0474 & 0.0409 & 0.0480 & 1.94 & $+48.0$ & 4/4 & 166 \\
1 & 4 & 4 & 0.0954 & 0.0465 & 0.0463 & 0.0530 & 2.05 & $+51.3$ & 4/4 & 191 \\
1 & 6 & 6 & 0.0960 & 0.0510 & 0.0523 & 0.0601 & 1.88 & $+46.8$ & 4/4 & 238 \\
2 & 1 & 2 & 0.0130 & 0.0132 & 0.0161 & 0.0109 & 0.98 & $-6.5$ & 1/4 & 224 \\
2 & 2 & 4 & 0.0127 & 0.0161 & 0.0160 & 0.0076 & 0.79 & $-29.0$ & 0/4 & 218 \\
2 & 4 & 8 & 0.0150 & 0.0243 & 0.0236 & 0.0103 & 0.62 & $-61.5$ & 0/4 & 209 \\
2 & 6 & 12 & 0.0246 & 0.0196 & 0.0195 & 0.0167 & 1.26 & $+20.4$ & 4/4 & 275 \\
4 & 1 & 4 & 0.0045 & 0.0253 & 0.0226 & 0.0079 & 0.18 & $-488.8$ & 0/4 & 212 \\
4 & 2 & 8 & 0.0082 & 0.0282 & 0.0279 & 0.0107 & 0.29 & $-229.8$ & 0/4 & 212 \\
4 & 4 & 16 & 0.0156 & 0.0397 & 0.0388 & 0.0171 & 0.39 & $-156.0$ & 0/4 & 165 \\
4 & 6 & 24 & 0.0249 & 0.0414 & 0.0414 & 0.0231 & 0.60 & $-69.1$ & 0/4 & 153 \\
\bottomrule
\end{tabular}
\par\smallskip\footnotesize $^{*}$\,Diagnostic oracle target (uses $F_{\rm ref}$); not deployable.
\end{table}


\begin{figure}[t]
  \centering
  \includegraphics[width=0.98\textwidth]{fig_wca_phase_01_heatmaps.png}
  \caption{Production WCA phase diagram at $M=100$: ABF final $\Ltwo(F)$ (left),
    mFR--ABF with the estimated target (middle), and the matched-seed improvement
    ratio $R=\Ltwo(F)_{\rm ABF}/\Ltwo(F)_{\rm mFR}$ (right). Values $R>1$ mean
    mFR improves over ABF. The gain localises in the high-error $\beta=1$ row,
    including the anchor cell; when ABF is already accurate, mFR is neutral to
    harmful.}
  \label{fig:wca_phase_heatmaps}
\end{figure}

The result is sharper than a monotone ``larger $\beta h$ is harder'' story.
The measured ABF error, not the nominal product $\beta h$, is the useful
starvation diagnostic. In the $\beta=1$ row ABF remains at
$\Ltwo(F)\approx0.09$--$0.10$, and mFR improves all matched seeds with
$+47$--$51\%$ median gain; the anchor cell $\beta=1,h=2$ improves from
$\WCAPDAnchorAbfErr$ to $\WCAPDAnchorFrErr$ ($R=\WCAPDAnchorR$). At
$\beta=2$ and $\beta=4$, ABF is often much more accurate (below
$\Ltwo(F)=0.005$ in the easiest accurate cell), and the same birth--death
configuration over-corrects: only $\WCAPDNWinCells/\WCAPDNPlaneCells$ cells have
positive gain, and the worst cell \WCAPDMinGainCell{} loses strongly. Thus the
phase diagram supports the conditional claim rather than an unconditional one:
mFR helps when the ABF free-energy estimator is sample-starved, but it can hurt
when ABF has already solved the marginal problem.

\begin{figure}[t]
  \centering
  \includegraphics[width=0.95\textwidth]{fig_wca_phase_03_transitions.png}
  \caption{Barrier-crossing counts for ABF and mFR--ABF (estimated target) in the
    same production cells. mFR slightly increases crossing counts in most cells, but
    the accuracy gain does not follow crossing count alone: unnecessary
    resampling can raise crossings while worsening $\Ltwo(F)$ when ABF is already
    accurate.}
  \label{fig:wca_phase_transitions}
\end{figure}

The mechanism diagnostics separate exploration from target-shape and genealogy
effects. The uniform target is close to the estimated target in this production run (mean absolute difference in matched gain
$\WCAPDMeanAbsUniMinusEst$ percentage points). The oracle target is different:
it improves most cells, which shows that birth--death can help beyond the
$\beta=1$ row if the target is accurate, but that target is diagnostic only. FR also increases
barrier crossings only modestly; the important pattern is that it reallocates
whole configurations in the cells where the ABF estimator is noisy. Genealogy does
not show catastrophic collapse in the production run: the minimum final ancestor ESS for
mFR--ABF is $\WCAPDMinFinalEss$ out of $1024$, the mean final ESS is
$\WCAPDMeanFinalEss$, and the largest single ancestor fraction observed is only
$\WCAPDMaxAncFrac$. The harmful cells therefore look like an unnecessary
resampling/variance tradeoff, not the severe genealogical collapse seen in the
aggressive-rate failure ladder.

\begin{figure}[t]
  \centering
  \includegraphics[width=0.95\textwidth]{fig_wca_phase_05_genealogy.png}
  \caption{Genealogical diagnostics for representative production regimes. The
    deployable mFR run keeps many ancestors alive in the helpful $\beta=1$ row,
    and the maximum ancestor fraction remains small; the production run does not exhibit
    the collapse associated with over-aggressive birth--death.}
  \label{fig:wca_phase_genealogy}
\end{figure}

\begin{figure}[t]
  \centering
  \includegraphics[width=0.95\textwidth]{fig_wca_phase_06_error_vs_time.png}
  \caption{Error-versus-time curves for representative cells. mFR separates from
    ABF after the FR onset in the sample-starved anchor regime, while in accurate
    high-$\beta$ regimes the same resampling is not beneficial.}
  \label{fig:wca_phase_time}
\end{figure}

This production phase diagram therefore refines the WCA conclusion. There is an
``easy/accurate'' region where ABF already works well and mFR should be reduced or
disabled, a sample-starved region (the $\beta=1$ row, including the original WCA
anchor) where gentle mFR gives a large free-energy gain, and intermediate cells
where the answer depends on the online target. The targeted mass slice at
$(\beta,h)=(2,4)$ is non-monotone: mFR helps at $M=49$ ($+24.4\%$), but hurts at
$M=100$ and $M=196$, while the oracle target helps the two larger systems. Thus
physical size alone is not the difficulty variable; the relevant diagnostic is
whether the practical online target is correcting a genuinely starved ABF
estimator without adding unnecessary birth--death noise. This is distinct from the
separate over-aggressive failure mode identified by the rate ladder, where
diversity collapse dominates.

\subsection{Case finding}
At the original condensed-phase WCA anchor, ABF+FR with an online estimated target
reliably and substantially improves free-energy accuracy: $+\WCAtunedGainPct\%$ in
median final $\Ltwo(F)$ at the tuned operating point, winning
$\WCAtunedWin/10$ seeds, with the gain concentrated in the free energy (not the
mean force). The mechanism is balanced-resampling variance reduction after
warm-up; at the anchor, target shape matters little
($\text{uniform}\approx\text{estimated}\approx\text{oracle}$, and the estimated
target already reaches the oracle's accuracy). The method has a clear safe regime,
\texttt{fr\_rate}$\approx0.1$--$0.2$, and a real failure boundary: too-aggressive
birth--death reverses the benefit as diversity collapses. The new phase diagram
qualifies the claim: mFR is most useful where the ABF estimator is visibly
sample-starved (the $\beta=1$ WCA row in the production phase diagram), but when ABF already has a low
free-energy error the same birth--death correction can be unnecessary or harmful.
Thus the practical recommendation is conditional, not automatic: use gentle,
diversity-preserving mFR in starved WCA regimes, and monitor ABF error,
transition support, FR event fractions, and genealogy before increasing the
birth--death strength.
